\chapter{Automate the Cross-Jail Integration}

The manual browser check proved the feature once. We now want a repeatable
test that boots the store in its test environment, dispatches a request through
its Rails stack, calls the real catalog jail, and asserts the rendered result.

This is deliberately a \emph{live integration test}. It crosses from the store
jail to the catalog service and its database, but it dispatches the incoming
store request in-process through Rails. A later external runner will enter
through the web jail and complete the full edge-to-edge path.

\section{Exercise 14A: Replace the generated controller test}

Keep the catalog server running. The store Puma process is not required for
this focused test because Rails' integration-test harness dispatches the store
request in-process.

In the store jail, open the controller test generated in the previous chapter:

\begin{lstlisting}[language=sh]
su -l store
cd /usr/local/www/store
vi test/controllers/products_controller_test.rb
\end{lstlisting}

Replace its contents with:

\begin{lstlisting}[language=Ruby]
require "test_helper"

class ProductsControllerTest < ActionDispatch::IntegrationTest
  test "renders products from the live catalog" do
    skip "set LIVE_CATALOG=1 to run" unless
      ENV["LIVE_CATALOG"] == "1"

    get root_url

    assert_response :success
    assert_select "h1", "Products"
    assert_select "li", text: /FreeBSD Handbook/
    assert_select "li", text: /BSD-001/
    assert_select "li", text: /\$49\.99/
  end
end
\end{lstlisting}

The environment-variable gate keeps an ordinary unit-test run deterministic
when the catalog service is intentionally stopped. Enabling the gate makes the
external dependency part of the test contract.

\section{Exercise 14B: Run the live test}

Run only this test with the live-catalog gate enabled:

\begin{lstlisting}[language=sh]
LIVE_CATALOG=1 bin/rails test \
  test/controllers/products_controller_test.rb
\end{lstlisting}

Expect one passing test. Its assertions prove that the store received a
successful response from the live catalog and converted the returned product
data into the expected HTML.

Now run the ordinary suite without opting into external services:

\begin{lstlisting}[language=sh]
bin/rails test
\end{lstlisting}

The live test should be reported as skipped rather than failing because a
dependency was not requested. This separation lets fast local tests and live
cross-jail tests coexist in the same application.

\section{Checkpoint}

Stop after the gated test passes and the ordinary suite skips it cleanly. We
now have a repeatable store-to-catalog integration check. The remaining
infrastructure work is to place the web jail in front of the store and run a
black-box test from outside both Rails applications.
